%% pcss_joint_definition.tex — PCSS Theory: L4 Joint Security Predicate
%% Paper F: "Physically-Constrained Semantic Security Theory for IoT Mesh Routing"
%% Issue #89 (F10): Joint PCSS-secure(P,E) definition integrating SBC + Physical Feasibility
%%
%% Depends on:
%%   definitions.tex         — Definitions 1–6 (Node State, D_r, Sem, Cl, DC, SBC)
%%   physical_feasibility.tex — Definitions 7–9 (E, C_auth, Phi), Theorem 2 (rho*)
%%   theorem1_proof.tex      — Theorem 1 (SBC -> next_hop ≠ A), 2-node
%%   theorem1_multihop.tex   — Theorem 1' (multi-hop), Corollary 2
%%
%% Required preamble packages:
%%   \usepackage{amsthm, amsmath, amssymb, mathtools}
%%
%% Theorem environments (declare in preamble if not already present):
%%   \theoremstyle{definition}
%%   \newtheorem{definition}{Definition}[section]
%%   \newtheorem{remark}[definition]{Remark}
%%   \newtheorem{example}[definition]{Example}
%%   \newtheorem{proposition}[definition]{Proposition}
%%   \theoremstyle{plain}
%%   \newtheorem{theorem}{Theorem}[section]
%%   \newtheorem{corollary}[theorem]{Corollary}

% -----------------------------------------------------------------------
\subsection{L4 — Joint PCSS Security Predicate}
\label{subsec:l4-pcss-joint}
% -----------------------------------------------------------------------

The definitions and theorems of Sections~\ref{subsec:l1-semantics}–\ref{subsec:l3-feasibility}
establish two orthogonal conditions that a protocol must satisfy for its routing tables to be
immune from adversarial manipulation in a concrete IoT deployment.  The
\emph{Semantic Binding Condition} (Definition~\ref{def:sbc}) captures the
cryptographic dimension: every decision-critical message must carry an
unforgeable, fresh authenticator.  The \emph{Physical Feasibility
Predicate} (Definition~\ref{def:feasibility}) captures the resource
dimension: the authentication rate mandated by the protocol must be
sustainably achievable within the duty-cycle and battery constraints of
the target hardware.

Neither condition alone is sufficient for a deployment guarantee:
\begin{itemize}
  \item $\mathrm{SBC}(P)$ holds but $\Phi(P,E)=0$: the protocol is
        cryptographically correct in the abstract model yet cannot be
        executed at the required rate on the actual hardware.  The
        implementation must silently reduce the authentication rate,
        re-opening the replay attack surface that SBC was designed to close.

  \item $\Phi(P,E)=1$ but $\mathrm{SBC}(P)$ fails: the hardware can
        sustain the required rate, but the protocol does not authenticate
        its decision-critical messages, so adversarial injection or
        field-manipulation is trivially possible.
\end{itemize}

We therefore combine both conditions into a single predicate, which we
call \emph{Physically-Constrained Semantic Security} (PCSS).

% -----------------------------------------------------------------------
\begin{definition}[PCSS-secure]
\label{def:pcss-secure}
Let $P$ be a routing protocol and $E$ an environment model
(Definition~\ref{def:environment}).
$P$ is \emph{PCSS-secure in environment $E$}, written
$\mathrm{PCSS\text{-}secure}(P,E)$, if and only if both of the following
hold simultaneously:
\begin{enumerate}[label=\textbf{PC\arabic*}., ref=PC\arabic*]
  \item \label{pc:sbc}
        \textbf{Semantic Binding.}\quad
        $P$ satisfies the Semantic Binding Condition:
        \[
          \mathrm{SBC}(P).
        \]
        That is, every message in $\mathrm{DC}(P)$ is authenticated,
        fresh, and has no unauthenticated equivalent
        (Definition~\ref{def:sbc}, conditions \ref{sbc:auth}–\ref{sbc:closure}).

  \item \label{pc:phi}
        \textbf{Physical Feasibility.}\quad
        The required authentication rate is within the hardware
        resource budget:
        \[
          \Phi(P, E) = 1,
        \]
        i.e., $\rho_{\mathrm{DC}}(P) \le \rho^*(E)$
        (Definition~\ref{def:feasibility}, Theorem~\ref{thm:feasibility-bound}).
\end{enumerate}
Formally:
\[
  \mathrm{PCSS\text{-}secure}(P, E)
  \;\;\Longleftrightarrow\;\;
  \mathrm{SBC}(P) \;\wedge\; \Phi(P, E) = 1.
\]
\end{definition}

\begin{remark}[$\mathrm{PCSS\text{-}secure}$ as a conjunction of independent layers]
\label{rem:pcss-independence}
The two conditions are \emph{logically independent}: satisfying one does
not imply the other.  SBC is a purely syntactic property of the protocol
specification (does every DC message carry an authenticator?), whereas
$\Phi$ is a numeric inequality over hardware constants and the
protocol-mandated transmission rate.  A protocol can be redesigned to
satisfy SBC without altering $\rho_{\mathrm{DC}}$, and vice versa.
The independence is not merely formal: in practice, adding HMAC to
LoRaMesher Hello messages (fixing SBC) does \emph{not} automatically
change the Hello interval (leaving $\Phi$ potentially violated).  PCSS
forces both issues to be addressed simultaneously.
\end{remark}

% -----------------------------------------------------------------------
\begin{theorem}[PCSS-secure is sufficient for attacker exclusion]
\label{thm:pcss-sufficiency}
Let $\mathrm{Assumption~1}$ (EUF-CMA security of the MAC scheme)
hold~\cite{goldwasser1988digital}.
If $\mathrm{PCSS\text{-}secure}(P, E)$ and attacker $\mathcal{A}$ is not
in $\mathrm{AuthorizedSet}(P)$, then for every execution trace $t$ and
every node $v_i \in V$:
\[
  \mathrm{next\_hop}_{v_i}(t) \neq \mathcal{A}.
\]
\end{theorem}

\begin{proof}
$\mathrm{PCSS\text{-}secure}(P, E)$ implies $\mathrm{SBC}(P)$ by
condition~\ref{pc:sbc} of Definition~\ref{def:pcss-secure}.  Apply
Theorem~\ref{thm:sbc-multihop} (Theorem~1', multi-hop extension)
directly: under $\mathrm{SBC}(P)$ and Assumption~1, no node $v_i \in V$
installs $\mathcal{A}$ as a next-hop entry for any destination in any
execution trace, provided $\mathcal{A} \notin \mathrm{AuthorizedSet}(P)$.

The physical feasibility condition $\Phi(P, E) = 1$
(condition~\ref{pc:phi}) is required to guarantee that $\mathrm{SBC}(P)$
holds \emph{in the concrete implementation}: since
$\rho_{\mathrm{DC}}(P) \le \rho^*(E)$, the node can authenticate every
$\mathrm{DC}(P)$ message at the protocol-mandated rate without
duty-cycle violation or battery exhaustion.  No authenticated message is
withheld, so the freshness condition (SBC2) is never silently weakened by
the implementation.  Therefore $\mathrm{SBC}(P)$ is not merely nominal
but operationally enforced, and the hypothesis of Theorem~\ref{thm:sbc-multihop}
is satisfied without qualification.
\hfill$\square$
\end{proof}

\begin{remark}[Why $\Phi$ appears in the proof]
\label{rem:phi-in-proof}
The proof step invoking $\Phi(P,E)=1$ is not a formality.  If instead
$\Phi(P,E) = 0$, the implementation must reduce the effective
authentication rate $\rho' < \rho_{\mathrm{DC}}(P)$; the freshness
condition SBC2 no longer holds at every message, and the conclusion of
Theorem~\ref{thm:sbc-multihop} no longer applies.  The attacker can
replay a stale but validly-signed $\mathrm{DC}(P)$ message captured
within the resulting replay window $\Delta t = 1/\rho'$.  This is
precisely the attack captured by the Proposition below.
\end{remark}

% -----------------------------------------------------------------------
\begin{proposition}[Replay window under $\Phi=0$]
\label{prop:replay-window}
Suppose $\mathrm{SBC}(P)$ holds in the abstract model (all $\mathrm{DC}(P)$
messages are specified to carry a valid HMAC and freshness token), but
$\Phi(P, E) = 0$, i.e., $\rho_{\mathrm{DC}}(P) > \rho^*(E)$.
Let $\rho' \le \rho^*(E)$ be the reduced authentication rate that the
implementation must use to remain within hardware constraints.

Then:
\begin{enumerate}[label=(\roman*)]
  \item \textbf{Replay window.}  Between consecutive authenticated
        transmissions of a $\mathrm{DC}(P)$ message for a given
        destination, there is a window of duration
        \[
          \Delta t \;=\; \frac{1}{\rho'} \;\ge\; \frac{1}{\rho^*(E)}
          \quad\text{seconds}
        \]
        during which any previously captured, validly-signed
        $\mathrm{DC}(P)$ message can be re-injected.

  \item \textbf{Attack feasibility.}  Within this window, attacker
        $\mathcal{A}$ (who has recorded a valid
        $m_{\mathrm{DC}} \in \mathrm{DC}(P)$ with \texttt{metric}$=0$
        from a prior round) can replay $m_{\mathrm{DC}}$.  Since the
        freshness check of SBC2 is momentarily unenforceable (no fresher
        authenticated update has been received), any node $v_i$ that
        accepts the replayed message may install
        $\mathcal{RT}_{v_i}(v_d) \leftarrow \mathcal{A}$ — a
        \emph{black-hole attack}.

  \item \textbf{Dolev-Yao blindness.}  In the classical Dolev-Yao
        model the replay attack is \emph{impossible by construction}:
        authenticated messages cannot be replayed because the model
        assumes instantaneous and universally enforced freshness checks.
        $\mathrm{PCSS\text{-}secure}$ is the minimal extension that
        makes the attack \emph{predictable}: the gap $\Phi(P,E)=0$ is a
        quantitative warning that the designer must either reduce
        $\rho_{\mathrm{DC}}(P)$ or increase $\rho^*(E)$ (e.g., by
        upgrading the battery or duty-cycle budget).
\end{enumerate}

\noindent Formally, when $\Phi(P,E)=0$ the security guarantee degrades
from the unconditional
\[
  \forall t,\; \forall v_i \in V:\;
  \mathrm{next\_hop}_{v_i}(t) \neq \mathcal{A}
  \quad (\text{Theorem~\ref{thm:pcss-sufficiency}})
\]
to the weaker, window-parameterised statement:
\[
  \Pr\!\bigl[\,
    \mathrm{next\_hop}_{v_i}(t) = \mathcal{A}
    \;\mid\;
    \mathcal{A} \text{ replays within } \Delta t
  \,\bigr]
  \;\ge\; p_{\mathrm{replay}}(\Delta t, n),
\]
where $p_{\mathrm{replay}}(\Delta t, n) > 0$ depends on the topology
and attacker radio range, and is \emph{not} negligible for typical
LoRa single-hop scenarios.
\end{proposition}

\begin{remark}[PCSS as the strongest routing-security guarantee]
\label{rem:three-tier}
The PCSS framework induces a strict three-tier security hierarchy:

\medskip
\noindent\textbf{Tier 1 — Dolev-Yao (DY) security.}\quad
The attacker cannot forge authenticated messages (computational hardness
assumption on the MAC).  DY guarantees that an adversary with unbounded
communication but bounded computation cannot produce a valid tag for a
message it did not receive.  \emph{DY does not predict degradation under
resource constraints}: it assumes the honest protocol runs at whatever
rate the specification mandates, without querying whether the hardware
can sustain that rate.

\medskip
\noindent\textbf{Tier 2 — SBC alone.}\quad
Every message in $\mathrm{DC}(P)$ is fully authenticated and fresh.  SBC
is strictly stronger than standard DY message authentication: it
additionally requires that \emph{every decision-critical field} is bound
by the authenticator (condition~\ref{sbc:closure}) and that freshness
tokens are universally enforced (condition~\ref{sbc:fresh}).  However,
SBC is a property of the \emph{protocol specification}; it does not check
whether the specification is realizable within the hardware resource
envelope of the target deployment.

\medskip
\noindent\textbf{Tier 3 — PCSS-secure.}\quad
$\mathrm{SBC}(P) \wedge \Phi(P,E)=1$.  This is the only tier that
provides a \emph{full deployment guarantee}: authentication is both
cryptographically sound (SBC) and physically sustainable ($\Phi=1$) in
the target environment.  Theorem~\ref{thm:pcss-sufficiency} shows that
Tier~3 implies the attacker-exclusion property of Tier~2, which in turn
subsumes Tier~1.  The converse does not hold: Tier~1 $\not\Rightarrow$
Tier~2 $\not\Rightarrow$ Tier~3.

\medskip
The hierarchy is summarised in Table~\ref{tab:security-hierarchy}.
\end{remark}

\begin{table}[h]
\centering
\caption{Three-tier security hierarchy in the PCSS framework.
         A checkmark (\checkmark) indicates that the tier provides the
         listed guarantee; a cross ($\times$) indicates it does not.}
\label{tab:security-hierarchy}
\small
\begin{tabular}{lccc}
  \hline
  \textbf{Guarantee} & \textbf{Tier 1 (DY)} & \textbf{Tier 2 (SBC)} & \textbf{Tier 3 (PCSS)} \\
  \hline
  Forgery impossible (EUF-CMA)
    & \checkmark & \checkmark & \checkmark \\
  All DC fields authenticated (SBC1+SBC3)
    & $\times$   & \checkmark & \checkmark \\
  Replay unconditionally prevented (SBC2)
    & $\times$   & \checkmark & \checkmark \\
  Authentication sustainable on hardware ($\Phi=1$)
    & $\times$   & $\times$   & \checkmark \\
  Attacker exclusion in concrete deployment
    & $\times$   & $\times$   & \checkmark \\
  Replay window predicted / quantified
    & $\times$   & $\times$   & \checkmark \\
  \hline
\end{tabular}
\end{table}

% -----------------------------------------------------------------------
\begin{corollary}[LoRaMesher deployment is not PCSS-secure]
\label{cor:loramesher-not-pcss}
Let $P_{\mathrm{LM}}$ be the default LoRaMesher protocol and
$E_{\mathrm{EoRa}}$ the EoRa-PI environment
(Example~\ref{ex:eora-pi}).  Then
$\mathrm{PCSS\text{-}secure}(P_{\mathrm{LM}}, E_{\mathrm{EoRa}})$ is
\emph{false}, by two independent violations:

\begin{enumerate}[label=(\roman*)]
  \item \textbf{SBC violation.}
        As shown in Example~\ref{ex:loramesher},
        $\mathrm{DC}(P_{\mathrm{LM}}) = \{m_{\mathrm{Hello}}\}$ and
        $m_{\mathrm{Hello}} \notin \mathrm{Auth}(P_{\mathrm{LM}})$, so
        condition~\ref{sbc:auth} of Definition~\ref{def:sbc} fails.
        Hence $\mathrm{SBC}(P_{\mathrm{LM}}) = \mathbf{false}$,
        directly violating condition~\ref{pc:sbc} of
        Definition~\ref{def:pcss-secure}.

  \item \textbf{Physical infeasibility.}
        Even if $P_{\mathrm{LM}}$ is patched to add HMAC-SHA256 to
        every Hello message (so that SBC1--SBC3 are satisfied),
        the default Hello interval of 30~s yields
        $\rho_{\mathrm{DC}}(P_{\mathrm{LM}}) = 1/30 \approx
        3.33 \times 10^{-2}\ \mathrm{pkt/s}$.
        Since $\rho^*(E_{\mathrm{EoRa}}) = 1.69 \times 10^{-3}\ \mathrm{pkt/s}$
        (Example~\ref{ex:eora-pi}), we have
        \[
          \rho_{\mathrm{DC}}(P_{\mathrm{LM}})
          \;=\; 3.33 \times 10^{-2}
          \;\approx\; 20\,\rho^*(E_{\mathrm{EoRa}}),
        \]
        so $\Phi(P_{\mathrm{LM}}, E_{\mathrm{EoRa}}) = 0$, violating
        condition~\ref{pc:phi} of Definition~\ref{def:pcss-secure}.
\end{enumerate}

Therefore $\mathrm{PCSS\text{-}secure}(P_{\mathrm{LM}}, E_{\mathrm{EoRa}})
= \mathbf{false}$ regardless of which violation is considered.  By
Proposition~\ref{prop:replay-window} the effective replay window in the
physically-constrained implementation is at least $\Delta t \ge 593\ \mathrm{s}$,
during which a stale Hello with \texttt{metric}$=0$ can install the
attacker as a next-hop.  This is confirmed by the ns-3 simulations of
Phase~4 (ns-3 simulation results from Paper~B).
\end{corollary}

% -----------------------------------------------------------------------
\begin{corollary}[Patched LoRaMesher at reduced Hello rate is PCSS-secure]
\label{cor:loramesher-patched-pcss}
Let $P_{\mathrm{LM}}^{+}$ denote the patched LoRaMesher protocol
obtained by (a)~appending an HMAC-SHA256 tag to every Hello message
(satisfying SBC1 and SBC3) and (b)~including a per-destination
MetricVersion counter in each Hello message (satisfying SBC2), and
let the Hello transmission rate be reduced to
$\rho \le \rho^*(E_{\mathrm{EoRa}}) = 1.69 \times 10^{-3}\ \mathrm{pkt/s}$
(i.e., at most one Hello per 593~s $\approx 9.9\ \mathrm{min}$).
Then $\mathrm{PCSS\text{-}secure}(P_{\mathrm{LM}}^{+}, E_{\mathrm{EoRa}}) = \mathbf{true}$.
\end{corollary}

\begin{proof}
We verify both conditions of Definition~\ref{def:pcss-secure}:

\medskip
\noindent\textbf{Condition~\ref{pc:sbc} — SBC.}\quad
By construction:
\begin{itemize}
  \item \emph{SBC1 (authenticity):} every Hello message carries an
        HMAC-SHA256 tag; under Assumption~1 (EUF-CMA), no polynomial-time
        adversary can forge a valid tag, so
        $\mathrm{DC}(P_{\mathrm{LM}}^{+}) \subseteq \mathrm{Auth}(P_{\mathrm{LM}}^{+})$.
  \item \emph{SBC2 (freshness):} the MetricVersion counter is
        monotonically non-decreasing per destination; any node rejects a
        Hello whose MetricVersion does not strictly advance the accepted
        counter.  No replay of a previously accepted Hello is possible.
  \item \emph{SBC3 (no unauthenticated equivalent):} the
        HMAC covers the entire Hello message including the
        \texttt{metric} and \texttt{next\_hop} fields; an adversary
        cannot alter these fields without invalidating the tag.
\end{itemize}
Hence $\mathrm{SBC}(P_{\mathrm{LM}}^{+}) = \mathbf{true}$.

\medskip
\noindent\textbf{Condition~\ref{pc:phi} — Physical Feasibility.}\quad
By assumption $\rho \le \rho^*(E_{\mathrm{EoRa}}) = 1.69 \times 10^{-3}\ \mathrm{pkt/s}$,
so $\rho_{\mathrm{DC}}(P_{\mathrm{LM}}^{+}) \le \rho^*(E_{\mathrm{EoRa}})$
and $\Phi(P_{\mathrm{LM}}^{+}, E_{\mathrm{EoRa}}) = 1$ by
Definition~\ref{def:feasibility}.

\medskip
Both conditions hold; Theorem~\ref{thm:pcss-sufficiency} therefore
guarantees that no unauthorised node is ever installed as a next-hop in
any execution trace.
\hfill$\square$
\end{proof}

\begin{remark}[Trade-off: routing convergence vs.\ security]
\label{rem:convergence-tradeoff}
Corollary~\ref{cor:loramesher-patched-pcss} establishes PCSS-security
at the cost of a significantly increased Hello interval
($\ge 593\ \mathrm{s}$ vs.\ the default $30\ \mathrm{s}$).  For a
distance-vector protocol on an $n$-node network, routing convergence
after a topology change requires $O(n)$ DV rounds, each lasting
$1/\rho$ seconds; the worst-case convergence time is
\[
  T_{\mathrm{conv}} \;=\; O\!\left(\frac{n}{\rho^*}\right)
  \;=\; O(n \times 593)\ \mathrm{seconds}.
\]
For the 6-node testbed ($n=6$) this gives
$T_{\mathrm{conv}} \approx 3558\ \mathrm{s} \approx 59\ \mathrm{min}$.
This trade-off is inherent to the physical constraints of the EoRa-PI
platform and cannot be resolved by protocol redesign alone.  Application
designers must either (a)~accept the increased convergence latency for
low-mobility or quasi-static deployments, (b)~upgrade to a
higher-capacity battery (which raises $\rho^*$ by increasing
$B_{\mathrm{mAh}}$), or (c)~reduce the control-message size (which
lowers $t_{\mathrm{tx}}$ and raises $\rho^*$).  The PCSS framework
makes the trade-off space explicit and quantitatively navigable.
\end{remark}
